We also ran the pipeline once on Google Cloud Dataproc, from ingestion to the Task 2 fit and validation. The 2024 test and every number reported elsewhere come from the local runs. Cluster \code{dsa5208-p2}, zone \code{asia-southeast1-c}, image \code{3.0.3-debian13}: 1 n2-standard-4 master and 2 n2-standard-4 workers, each with a boot disk of 80\,GB pd-balanced; Spark 4.1.2 on Java 21.0.12.1 and Python 3.12.13, master \code{yarn}; 2 executors with 3 cores and 8g each, driver memory 8g, 48 shuffle partitions, dynamic allocation off.

These three machines took the project's whole vCPU quota, so the cluster could not be larger. We sized the executors as Lecture 6 does. Each machine has 4 vCPUs and 16\,GiB of memory. Lecture 6 recommends three to five cores per executor and keeping some capacity on each worker for the cluster manager and the operating system, so we ran one executor per worker with 3 cores and left the remaining vCPU free. Executor memory, plus an overhead that Lecture 6 puts at about a tenth of it (more for PySpark), has to fit within what YARN offers on a worker, which is less than the machine's 16\,GiB; 8g fits with room to spare for the Python worker processes. We turned dynamic allocation off so that every job ran on exactly these executors. The raw CSVs were downloaded from data.gov.sg straight into Cloud Storage by a job on the cluster (\code{cloud/fetch\_raw\_on\_cluster.py}) rather than uploaded from our local machine. The run used two clusters of this configuration. The stages up to the cleaning-ledger check ran on the first (\code{cloud\_run/cluster\_part1.json}). A status call on our side hung, the cluster sat idle until its idle limit deleted it, and we resumed the remaining stages on a second cluster, which read the first cluster's outputs from the same Cloud Storage prefix.

Table~\ref{tab:cloudcmp} compares the wall time of each stage; the cloud stages took about 44 minutes in total. Every Dataproc job pays a fixed start-up cost (the trivial environment-check job took 29\,s), and on the short stages that cost dominates. Our local machine also has 10 cores against the cluster's 6 executor cores, and reads local disk rather than Cloud Storage. Cleaning, the longest data stage, is where the cluster catches up.

To check that the cluster computed the same results, \code{checks/cloud\_parity.py} compares its outputs with the local ones (22 tables). Counts that do not depend on the order of summation (rows, stations, every cleaning rule, coverage, slot and event counts, sample sizes) match exactly; across 34 floating-point columns the largest relative difference is $3.1\times10^{-13}$. The 4 threshold counts on summed rain (hours with at least 20 or 40\,mm, days with at least 50\,mm) differ by at most 4 per row. Most readings are multiples of 0.2\,mm, so many totals equal a threshold exactly, and a floating-point sum lands just above or below it depending on the order in which Spark adds partitions. \code{checks/threshold\_boundary.py} counts the totals within $10^{-6}$\,mm of each threshold, and all differences are within that count. The 5 approximate quantiles (Spark's \code{percentile\_approx}) differ by up to 1.3\,\%. On 2023 the Task 2 fit on the cluster chose the same models as the local fit; Model B's validation AUC-PR is 0.6677 on Dataproc and 0.6677 locally, and the largest AUC-PR difference over all candidates is 0.0023.

\begin{table}[H]\centering\small
\begin{tabular}{lrrr}
\toprule
Stage & Local (s) & Dataproc (s) & Dataproc / local \\
\midrule
ingest (CSV $\to$ Parquet) & 23 & 156 & 6.7 \\
ingest check & 19 & 45 & 2.3 \\
cleaning + ledger & 625 & 580 & 0.9 \\
ledger check & 14 & 48 & 3.4 \\
coverage & 4 & 51 & 13.4 \\
Task 1 aggregation & 63 & 143 & 2.3 \\
Task 2 features & 103 & 255 & 2.5 \\
Task 2 samples + baselines & 42 & 133 & 3.2 \\
Task 2 fit + validation & not recorded & 1251 & -- \\
\bottomrule
\end{tabular}

\caption{Wall time per stage in seconds, on the local machine and on Dataproc, and their ratio; a ratio below one means Dataproc was faster. \src{src/report\_macros.py from runs/*/env.json and cloud\_run/job\_times.json (Dataproc jobs API; driver start to exit, including the small Cloud Storage copies)}}\label{tab:cloudcmp}
\end{table}
